\chapter{Naučnik}

\editorialnote{Sol Vejntraub, Rejčel, Sarai, Merlinova bolest, vera, roditeljska ljubav i nemoguć izbor.}

Sol Vejntraub je dugo ćutao pre nego što je počeo.

Rejčel je spavala u improvizovanoj kolevci kraj njega.

Bila je toliko mala da je bilo teško povezati je sa odraslom ženom o kojoj je trebalo da govori.

Sol je spustio pogled na dete.

``Pre svega'', rekao je, ``morate da znate ko je ona bila.''

\bigskip

Sol i njegova supruga Sarai živeli su na Barnardovom Svetu.

Bio je to star, miran svet u Mreži, bez velikih gradova, neobičnih pejzaža ili događaja koji bi privlačili pažnju ostatka Hegemonije.

Njima je odgovarao.

Sol je predavao istoriju, filozofiju i etiku na koledžu Najtenhelzer.

Sarai je bila muzikološkinja i kritičarka.

Nisu bili bogati.

Nisu mnogo putovali.

Nisu želeli naročito uzbudljiv život.

Voleli su jedno drugo.

Pet godina pokušavali su da dobiju dete.

Onda se rodila Rejčel.

Sol je pamtio prve mesece kao mešavinu potpunog umora i sreće.

Često bi se noću probudio bez razloga, otišao do kolevke i samo gledao kako dete spava.

Ponekad bi tamo već zatekao Sarai.

Stajali bi jedno uz drugo, držali se za ruke i posmatrali svoju kćer.

Sol je godinama proučavao etiku, religiju i pitanje smisla.

Nijedna knjiga nije ga pripremila za jednostavnu činjenicu da čovek može nekoga da voli toliko da njegov sopstveni život odjednom prestane da bude najvažnija stvar na svetu.

Rejčel je rasla.

Bila je radoznala, tvrdoglava, pametna i potpuno sposobna da roditelje dovede do očajanja.

Sa pet godina ošišala je sve svoje lutke, a zatim i sebe.

Sa sedam je odlučila da siromašne porodice u okolini nemaju dovoljno hrane, ispraznila gotovo čitavu porodičnu ostavu i sve razdelila.

Sa deset se popela na drvo visoko nekoliko desetina metara samo zato što joj je jedan dečak rekao da ne može.

Grana je pukla.

Slomila je nogu, rebra i vilicu.

Kada je Sol stigao u bolnicu, Rejčel mu je kroz žice koje su joj držale vilicu objasnila da joj je do vrha nedostajalo svega nekoliko metara.

``Sledeći put ću uspeti.''

Takva je bila Rejčel.

\bigskip

Kao odrasla izabrala je arheologiju.

Niko se tome nije posebno iznenadio.

Još kao mala mogla je satima da kopa zemlju ispod porodične verande i donosi roditeljima svaku staru kovanicu ili odbačeni predmet koji pronađe, zahtevajući da joj objasne kome je pripadao i kakav je život ta osoba vodila.

Studirala je na Barnardovom Svetu, a zatim na Rajhs univerzitetu na Friholmu.

Specijalizovala se za predmete koji nisu bili ljudskog porekla.

To ju je, pre ili kasnije, moralo odvesti na Hiperion.

Kada im je prvi put rekla da ide tamo, Sarai je pokušala da zvuči opušteno.

``Zašto baš Hiperion?''

``Vremenske grobnice.''

Sol je odmah podigao pogled.

Čitao je o njima.

Drevne građevine.

Nepoznat graditelj.

Antientropijska polja.

Priče o Šrajku.

``Zar nisu opasne?''

Rejčel se nasmejala.

``Tata, niko ozbiljan više ne veruje u te priče.''

Sarai ju je pitala gde će boraviti.

Rejčel je objasnila da postoji istraživačka stanica i hotel kod Utvrđenja Hronos.

Sve je bilo organizovano.

Ekspediciju je vodio Melio Arundez, jedan od najvećih stručnjaka za Vremenske grobnice.

``Biće uzbudljivo'', rekla je.

``Ne previše'', odgovorila je Sarai.

Svi su se nasmejali.

Niko od njih nije znao koliko će dugo pamtiti taj razgovor.

\bigskip

Putovanje do Hiperiona značilo je godine vremenskog duga za ljude koji su ostali kod kuće.

Za Rejčel je put trajao samo nekoliko nedelja u fugi.

Za Sola i Sarai prošle su godine.

Za to vreme Sol je počeo da sanja.

U snu je hodao ogromnom građevinom.

Stubovi su nestajali u mraku visoko iznad njega.

Iz tame su ga posmatrale dve crvene oči.

A glas mu je govorio:

``Uzmi svoju jedinu kćer Rejčel, koju voliš. Odvedi je na Hiperion i prinesi je kao žrtvu.''

Sol je znao priču.

Abraham i Isak.

Bog traži od Abrahama da žrtvuje sina kako bi dokazao svoju veru.

Sol nije bio impresioniran.

``Ne.''

Glas je ponovio zahtev.

``Čuo sam te prvi put.''

Ponovo.

Budućnost čovečanstva zavisi od njegove poslušnosti.

Mora da povede Rejčel na Hiperion.

Mora da je preda.

Sol je u snu ugledao svoju kćer na kamenom oltaru.

U ruci je držao nož.

Bacio ga je.

``Ne.''

Probudio se uznemiren, ali je pokušao da san tretira kao ono što bi svaki racionalan čovek smatrao da jeste.

San.

Nekoliko meseci kasnije stigla je poruka sa Hiperiona.

Rejčel je doživela nesreću.

\bigskip

Ekspedicija je istraživala građevinu koju su ljudi nazvali Sfingom.

Bila je to jedna od Vremenskih grobnica.

Ogromna kamena struktura ispunjena hodnicima koji nisu vodili nikuda i menjali oblik i dimenzije bez očigledne logike.

Rejčel i Melio pokušavali su da mapiraju njenu unutrašnjost.

Jedne noći Rejčel je otišla sama.

Instrumenti su u početku pokazivali uobičajene vrednosti.

Onda se dogodilo nešto nemoguće.

Senzori su pokazali prostorije koje nekoliko trenutaka ranije nisu postojale.

Hodnici su se na instrumentima savijali sami u sebe.

Alarm za vremenske plime se aktivirao.

Svetla su se ugasila.

Ne samo električna svetla.

Sve.

Uređaji sa sopstvenim izvorima energije.

Njen komlog.

Čak i bioluminescentne oznake koje uopšte nisu zavisile od struje.

Rejčel je ostala u potpunom mraku.

Čula je nešto u hodniku.

Grebanje.

Metal o kamen.

Pokušala je da pronađe izlaz.

Tavanica se spustila.

Prolaz se zatvorio.

Nešto joj se približilo.

Poslednje čega se sećala bio je dodir nečega veoma hladnog oko ručnog zgloba.

I onda je vrisnula.

\bigskip

Kada su je pronašli, nije imala ozbiljne fizičke povrede.

Ali bila je bez svesti.

Prebačena je prema Mreži.

Sol i Sarai čekali su mesecima da stigne.

Kada su konačno došli do nje, lekari nisu znali kako da im objasne šta se dogodilo.

Rejčel se probudila.

Pogledala je roditelje.

``Mama? Tata? Šta vi radite ovde?''

Nije se sećala nesreće.

Nije se sećala transporta.

Mislila je da se još nalazi na Hiperionu.

Lekari su je pregledali ponovo.

I ponovo.

Konačno su dali ime onome što su videli.

Merlinova bolest.

Ime je bilo gotovo šala.

Problem je bio stvaran.

Rejčel je starila unazad.

Svakoga dana njeno telo bilo je približno jedan dan mlađe.

Ali to nije bilo najgore.

Svaki put kada bi zaspala, izgubila bi još jedan dan sećanja.

Ne dan koji je upravo proživela.

Nego poslednji dan sopstvenog života pre nesreće.

Sećanja su nestajala obrnutim redosledom.

Kao da neko briše njen život od kraja prema početku.

\bigskip

Rejčel je pokušala da se bori.

Snimala je poruke samoj sebi.

Svako jutro budila bi se zbunjena, uključila komlog i čula sopstveni glas.

``Dobro jutro, Rejč. Ne paniči.''

Zatim bi sama sebi objasnila šta se dogodilo.

Da je bila na Hiperionu.

Da ju je zahvatila vremenska plima.

Da njeno telo postaje mlađe.

Da će nakon sledećeg sna zaboraviti još jedan deo svog života.

Snimala je informacije o ljudima.

Roditeljima.

Meliju.

Ekspediciji.

O sopstvenoj vezi.

O tome šta voli.

Čega se plaši.

Šta je planirala.

Svako jutro bila je prisiljena da sluša biografiju osobe koja je izgledala kao ona, govorila kao ona i tvrdila da jeste ona.

Ali čijeg života se sve manje sećala.

Melio je pokušao da ostane uz nju.

Voleli su se.

Ali Rejčel je svaki dan gubila deo njihove veze.

Jednog jutra bi ga prepoznala kao čoveka kog voli.

Nekoliko nedelja kasnije kao čoveka kog je tek upoznala.

Posle toga kao potpunog stranca.

Fotografije nisu pomagale.

Snimci nisu pomagali.

Poruke koje je sama sebi ostavljala nisu mogle da stvore osećanje koje je izgubila.

Na kraju ga je napustila.

Ne zato što ga više nije volela.

Nego zato što više nije mogla da se seti da ga voli.

\bigskip

Vratila se roditeljima.

Sol ju je sačekao na terminalu.

Kada ga je ugledala, zastala je.

Za nju je od poslednjeg jasnog sećanja na oca prošlo mnogo manje vremena nego za njega.

On je ostario.

Kosa mu je nestala.

Lice se promenilo.

Rejčel je počela da plače.

``Izgledaš drugačije, tata.''

Sol je pokušao da se nasmeši.

``Prošlo je jedanaest godina, malena.''

Za nju nije.

To je bilo jedno od najokrutnijih pravila bolesti.

Svet je nastavljao napred.

Rejčel je išla unazad.

\bigskip

U početku je kod kuće bilo lakše.

Poznavala je kuću.

Ulice.

Koledž.

Roditelje.

Ali uskoro su se pojavili drugi problemi.

Za svoj ponovljeni dvadeset drugi rođendan želela je da vidi stare prijatelje.

Sol i Sarai su pristali.

Bila je to greška.

Ljudi kojih se Rejčel sećala kao svojih vršnjaka sada su bili mnogo stariji.

Neki su bili u braku.

Neki su imali decu.

Neki su se promenili toliko da ih je jedva prepoznala.

Neki su već bili mrtvi.

Rejčel je te večeri dugo sedela pred ostacima rođendanske torte.

Zatim je pogledala oca.

``Tata, nemoj više nikada da mi dozvoliš da uradim nešto ovakvo.''

Sledećeg jutra više se nije sećala proslave.

\bigskip

Jedne noći, dok je ponovo imala dvadeset jednu godinu, Rejčel je uzela sredstva da ne zaspi.

Želela je dovoljno vremena da pregleda sve snimke koje je ostavila samoj sebi.

Posle dva dana došla je do Solove sobe.

``Tata, možeš li da siđeš?''

Sedeli su u dnevnoj sobi.

Prvi i jedini put u životu pili su zajedno.

U početku su se smejali.

Pričali gluposti.

Onda je Rejčel postala ozbiljna.

``Pregledala sam sve.''

Sol nije odgovorio.

``Sve poruke. Medicinske izveštaje. Fotografije sa Hiperiona. Melija. Moje prijatelje.''

``I?''

``Ne verujem.''

Sol ju je pogledao.

``Znam da je istina'', rekla je. ``Ali ne osećam da je istina.''

Objasnila mu je.

Za nju je život koji je izgubila pripadao drugoj osobi.

Odrasla Rejčel koja je otišla na drugi svet.

Koja se zaljubila.

Koja je doživela nesreću.

Koja je mesecima pokušavala da spase svoje pamćenje.

Sve je to mogla intelektualno da prihvati.

Ali nije mogla svakog jutra ponovo da proživljava njenu patnju.

``Nemojte više da mi pričate.''

Sol je jedva razumeo.

``Šta?''

``Ti i mama. Nemojte da mi objašnjavate sve. Nemojte da puštate one poruke osim ako sama ne tražim.''

Počela je da plače.

``Rejčel koja je sve to doživela više nije ja. Zašto moram svakog dana ponovo da trpim njen bol?''

Sol ju je zagrlio.

``Razumem.''

I stvarno je razumeo.

Od tog dana prestali su da je prisiljavaju da pamti život koji joj je bolest već oduzela.

\bigskip

Rejčel je nastavila da postaje mlađa.

Dvadeset.

Osamnaest.

Šesnaest.

Ponovo je postala tinejdžerka.

Jednog jutra sišla je niz stepenice zabrinuta zbog ispita koji je polagala više od deset godina ranije.

Pitala je za prijatelja koji je odavno poginuo.

Pitala je za psa koji je već godinama bio mrtav.

Sol i Sarai naučili su da lažu.

Ne velike laži.

Male milosrdne laži.

``Danas nema škole.''

``Niki će se javiti kasnije.''

``Pas je kod veterinara.''

Nisu mogli svakog jutra da joj objasne da joj život nestaje.

Sol je sve vreme tražio lek.

Kontaktirao je lekare.

Fizičare.

Istraživačke institute.

Hegemoniju.

Niko nije imao odgovor.

Melio se vratio na Hiperion sa novom ekspedicijom.

Godinama je istraživao Sfingu.

Ništa.

Merlinova bolest nikada se nije ponovila ni kod jedne druge osobe.

Rejčel je bila jedina.

\bigskip

Kako je postajala mlađa, Sol je sve više razmišljao o snu.

Glas se i dalje vraćao.

``Odvedi Rejčel na Hiperion.''

Sol mu je sada odgovarao drugačije.

``Već si je uzeo!''

Tražio je objašnjenje.

Nije ga dobijao.

Počeo je ozbiljno da proučava priču o Abrahamu.

Čitavog života bavio se etikom.

Sada mu se staro pitanje više nije činilo akademskim.

Ako Bog naredi čoveku da ubije sopstveno dete, šta je moralno?

Poslušati Boga?

Ili zaštititi dete?

Razgovarao je sa rabinom.

Rabin mu je objasnio da Bog nije dozvolio Abrahamu da ubije Isaka.

Poenta je bila vera.

Poslušnost.

Sol se nije slagao.

``Ali Abraham je bio spreman da to uradi.''

Ako je moral moguć samo zato što čovek izvršava naređenje više sile, kakav je to moral?

Ako roditelj može da kaže:

Bog mi je naredio da ubijem svoje dete,

i time opravda postupak,

onda ljubav, odgovornost i slobodna volja ne znače ništa.

Sol je počeo da formuliše sopstveni odgovor.

Možda Abrahamova priča nije bila priča o tome kako čovek mora da posluša Boga.

Možda je došlo vreme da čovek kaže Bogu:

Ne.

\bigskip

Rejčel je postajala dete.

Sarai i Sol podvrgli su se tretmanima za podmlađivanje koliko su mogli, ne iz taštine, nego zato što nisu znali koliko dugo će morati da se brinu o njoj.

Rejčel je imala šesnaest.

Trinaest.

Deset.

Osam.

Melio se jednog dana vratio.

Rejčel je tada bila devojčica.

Sol ih je upoznao.

``Rejčel, ovo je dr Arundez.''

Ona se osmehnula.

``Rajhs univerzitet? Volela bih jednog dana tamo da studiram arheologiju.''

Melio nije mogao da govori.

Pred njim je stajala žena koju je voleo.

Ali ona još nije stigla do trenutka svog života u kojem će ga upoznati.

Kada je otišla u kuću, Melio je samo rekao:

``Izneverio sam je.''

``Nisi'', odgovorio je Sol.

Melio je otišao.

Sol ga više nikada nije video uživo.

\bigskip

Sol je pokušao i sa Crkvom Šrajka.

Otputovao je do njihovog najvećeg hrama.

Objasnio je biskupu šta se dogodilo Rejčel.

Kada je pomenuo da ju je vremenska plima zahvatila u Sfingi, sveštenikovo ponašanje se promenilo.

``Vaša kćer je odabrana.''

``Odabrana za šta?''

Biskup nije hteo da objasni.

Rekao je samo da joj niko ne može pomoći.

Da je istovremeno blagoslovena i prokleta.

Da ju je Bog bola izabrao da pati.

Sol je ustao.

``Moja kćer nije ničija žrtva.''

Izbacili su ga iz hrama.

Posle toga mu nijedan hram Šrajka nije želeo pomoći.

\bigskip

A onda je poginula Sarai.

Nije umrla od bolesti.

Nije bila žrtva Šrajka.

Nije postojao veliki kosmički razlog.

Jednostavno se dogodila nesreća.

Dvoje tinejdžera ukrali su vozilo.

Leteli su prebrzo.

Isključili sigurnosne sisteme da ih policija ne bi pratila.

Sudarili su se sa vozilom u kojem se nalazila Sarai.

Poginula je odmah.

Rejčel je tada imala četiri godine.

Odnosno, ponovo je imala četiri godine.

Sol joj je rekao.

Dete je plakalo.

Sutradan se probudilo.

``Gde je mama?''

Sol joj je ponovo rekao.

Rejčel je ponovo doživela majčinu smrt.

Trećeg jutra opet ga je pitala.

``Gde je mama?''

Sol je shvatio da ne može više.

Ne može da ubija Sarai pred njenom kćerkom svakog jutra.

``Mama je kod ujna Tete.''

``Hoću li je videti sutra?''

Sol je zastao.

``Da.''

Bio je to možda najteži odgovor koji je ikada dao.

\bigskip

Posle Sarainog pogreba Sol je otišao sam u pustinju.

Godinama je govorio da ne zna da li veruje u Boga.

Na pitanje Rejčel:

``Tata, veruješ li u Boga?''

uvek je odgovarao:

``Još čekam.''

Sada je prvi put shvatio da je čitavog života ipak razgovarao sa Njim.

Svađao se.

Optuživao.

Pitao zašto.

Možda je upravo to bila njegova vera.

Ne poslušnost.

Rasprava.

Odbijanje da prihvati patnju samo zato što mu neko kaže da ona ima smisla.

Sol je plakao za Sarai.

A sutradan je nastavio da brine o Rejčel.

\bigskip

Bolest je ubrzavala njen povratak u detinjstvo.

Stalni zubi su ispadali.

Mlečni zubi su se vraćali.

Zatim su i oni počeli da nestaju.

Prestala je da čita.

Prestala je samostalno da se oblači.

Izgubila je koordinaciju.

Jednog dana više nije mogla da hoda.

Najgori je bio jezik.

Reči su nestajale.

Sol je osećao kao da se sa svakom izgubljenom rečju ruši još jedan most između njih.

Jedne večeri, kada je Rejčel još mogla da govori, rekao joj je njihovu staru šalu.

``Vidimo se, aligatoru.''

Ona ga je zbunjeno pogledala.

Više se nije sećala odgovora.

Sol joj ga je ponovo objasnio.

``Kasnije, krokodilče.''

Rejčel se nasmejala.

``Kasnije, 'kodilče.''

Sutradan je ponovo zaboravila.

Njena poslednja jasna reč bila je:

``Mama.''

\bigskip

Na kraju je ponovo bila beba.

Sol ju je hranio.

Menjao joj pelene.

Nosio je uz grudi.

I čudna stvar, rekao je ostalima, bila je da je uprkos svemu voleo te trenutke.

Nije voleo bolest.

Nije prihvatao ono što se događa.

Ali voleo je svoju kćer.

I kada bi mu zaspala na grudima, sa glavom uz njegov obraz, na nekoliko minuta bolest bi prestajala da bude najvažnija stvar na svetu.

Postojali su samo otac i dete.

Kao pre mnogo decenija.

Kada se rodila.

Samo što je sada znao da sat ide prema nuli.

\bigskip

Počeo je da moli za dozvolu da ode na Hiperion.

Crkva Šrajka ga je odbijala.

Hegemonija je odugovlačila.

Sol je postao javan.

Pojavljivao se u medijima sa Rejčel u naručju.

``Ne znam da li će joj Hiperion pomoći'', govorio je. ``Ali tamo je sve počelo. Tražim samo priliku da pokušam.''

Pritisak javnosti je rastao.

A onda je stigla vest.

Crkva je odobrila poslednje hodočašće.

Sol Vejntraub i njegova kćer bili su među izabranima.

Rejčel je tada imala samo nekoliko nedelja života pre nego što bi se, prema svim proračunima, vratila preko trenutka sopstvenog rođenja.

Put do Hiperiona trajao je gotovo isto toliko.

Nije postojalo vreme za razmišljanje.

Sol je krenuo.

\bigskip

Na brodu prema Hiperionu ponovo je sanjao.

Hodao je istom mračnom građevinom.

Ovoga puta držao je bebu Rejčel u naručju.

Ponovo se začuo glas.

``Uzmi svoju jedinu kćer, koju voliš, i prinesi je kao žrtvu.''

Sol nije odmah odgovorio.

Pogledao je Rejčel.

Njene oči bile su otvorene.

Imao je osećaj da ga razume.

Tada je podigao glavu.

``Ne.''

Glas je čekao.

``Nema više žrtvovanja dece.''

Sol je držao Rejčel čvršće.

``Nema više slepe poslušnosti. Nema više Boga koji od roditelja zahteva da mu predaju dete kako bi dokazali da Ga vole.''

Tama je ćutala.

Sol je nastavio.

``Ako postoji novi savez između Boga i čoveka, onda i Bog mora nešto da nauči.''

Stubovi oko njega počeli su da se tresu.

``Abraham je bio spreman da žrtvuje sina zbog Boga. Ja neću.''

Glas se nije čuo.

``Ako želiš ljubav'', viknuo je Sol, ``onda nauči šta znači biti roditelj.''

Crvene oči su bledele.

``Neću ti predati Rejčel kao žrtvu.''

Privukao ju je uz grudi.

``Ako želiš da razumeš čovečanstvo, pridruži nam se kao otac. Nemoj čekati kao bog koji zahteva krv.''

Tlo se zatreslo.

Tama se zatvorila oko njih.

Sol se probudio.

Rejčel je bila u njegovom naručju.

Pogledala ga je.

I osmehnula se.

Sol je za trenutak prestao da govori.

Pogledao je bebu koja je sada spavala kraj njega.

``To je bio njen poslednji osmeh'', rekao je.

Zastao je.

``I njen prvi.''

\bigskip

Niko u kabini nije progovorio.

Bron Lamija je obrisala oči i pogledala u stranu.

Čak je i Martin Silenus ćutao.

Konzul je gledao dete.

Sada je potpuno razumeo zbog čega je Sol Vejntraub išao prema Šrajku.

Ne zbog vere.

Ne zbog poslušnosti.

Čak ni zato što je verovao da će Šrajk ispuniti želju.

Išao je zato što su sve druge mogućnosti nestale.

Hojt je konačno progovorio.

``Ako stignemo do Vremenskih grobnica... šta ćete uraditi?''

Sol je dugo gledao svoju kćer.

``Ne znam.''

``Ali san vam je rekao da je predate.''

``Da.''

``Šrajku?''

Sol je klimnuo.

``Mislite li da je glas bio Bog?''

Sol se blago osmehnuo.

``Čitavog života sam čekao dokaz da Bog postoji.''

Pogledao je u tamni prozor vetrokola.

``Sada se više plašim mogućnosti da postoji.''

Hojt nije odgovorio.

``Ali jedno znam'', nastavio je Sol. ``Ako bilo koji bog zahteva od mene da ubijem dete kako bih dokazao svoju ljubav prema njemu, onda taj bog nije zaslužio moju poslušnost.''

``A ako je predaja Rejčel jedini način da je spasete?'' upitala je Lamija.

Sol je spustio glavu.

Tu se nalazilo pitanje od kojeg nije mogao da pobegne.

Ako odbije, Rejčel će nestati.

Ako posluša, možda je predaje upravo onome što ju je osudilo.

Roditeljska ljubav govorila mu je da je zaštiti.

Nada ga je terala da je odnese do kraja puta.

A razum mu nije nudio nikakav odgovor.

``Zato sam ovde'', rekao je.

Pružio je ruku i nežno dodirnuo bebinu šaku.

Njeni prsti automatski su se zatvorili oko njegovog.

``Ne znam šta ću uraditi kada dođe trenutak.''

Pogledao je ostale.

``Znam samo da do tog trenutka moram da je donesem živu.''

Vetrokola su nastavila da tutnje kroz tamu.

Negde ispred njih nalazile su se Vremenske grobnice.

A u Solovom naručju spavala je njegova kćer, svakog minuta sve bliža početku sopstvenog života.